\documentclass[11pt]{article}
\usepackage{amsmath,amssymb,amsthm}
\setlength{\emergencystretch}{2em}
\newtheorem{theorem}{Theorem}
\newtheorem{corollary}[theorem]{Corollary}
\newtheorem{proposition}[theorem]{Proposition}
\newcommand{\End}{\operatorname{End}}
\newcommand{\im}{\operatorname{im}}
\newcommand{\Fix}{\operatorname{Fix}}
\begin{document}
\title{Canonical fixed vectors in characteristic two}
\author{Human Compression}
\date{}
\maketitle

We give a conditional nongeneration criterion requiring neither a
centralizer order nor a rank-one operator.  The construction uses elementary
properties of Frobenius, adjoints, and orthogonal complements; no claim of
originality is made for these algebraic ingredients.  The proposed use is
to extract a fixed vector from an algebra generated by two specified group
elements, and to find such certificates by linear equations.  There is no
claim here that a retained finite-group instance satisfies the hypotheses.

Let $k$ be a perfect field of characteristic two.  Thus $t\mapsto t^2$ is
an automorphism of $k$, and we write $\sqrt{t}$ for its inverse.  Finite
fields of characteristic two satisfy this hypothesis.  Let $V$ be a
finite-dimensional space with a nondegenerate symmetric bilinear form
$\beta$.  We use column vectors throughout.

\begin{theorem}\label{thm:canonical}
For each self-adjoint $P\in\End_k(V)$ there is a unique vector $w(P)$
such that
\[
             \beta(w(P),x)^2=\beta(Px,x)\qquad(x\in V).
\]
This vector lies in $\im P$.  For every linear isometry $h$ of $\beta$,
\[
                  w(hPh^{-1})=h\,w(P).
\]
Consequently, if $H$ is a group of linear isometries commuting with $P$,
then $H$ fixes $w(P)$.
\end{theorem}

\begin{proof}
Set $q_P(x)=\beta(Px,x)$.  For $x,y\in V$, the two cross terms in
$q_P(x+y)$ are equal: self-adjointness and symmetry give
$\beta(Px,y)=\beta(x,Py)=\beta(Py,x)$.  Their sum is zero in
characteristic two.  Hence
\[
 q_P(x+y)=q_P(x)+q_P(y),\qquad q_P(t x)=t^2q_P(x).
\]
Since inverse Frobenius is an automorphism, the function
$\ell_P(x)=\sqrt{q_P(x)}$ is a $k$-linear functional.  Nondegeneracy
gives a unique vector $w(P)$ with $\beta(w(P),x)=\ell_P(x)$ for all $x$.

If $x\in\ker P$, then $q_P(x)=0$, so $\beta(w(P),x)=0$.
Thus $w(P)\in(\ker P)^\perp$.  Self-adjointness implies
$\im P\subseteq(\ker P)^\perp$: for $z\in\ker P$ one has
$\beta(Py,z)=\beta(y,Pz)=0$.  Both spaces have dimension
$\dim V-\dim\ker P$, so they are equal.  This proves $w(P)\in\im P$.

If $h$ is an isometry, then
\[
 q_{hPh^{-1}}(x)=\beta(hPh^{-1}x,x)
                =\beta(Ph^{-1}x,h^{-1}x)=q_P(h^{-1}x).
\]
Taking square roots gives
$\ell_{hPh^{-1}}(x)=\ell_P(h^{-1}x)=\beta(hw(P),x)$.
Uniqueness proves the equivariance formula.  If $hP=Ph$, this formula
says precisely that $h\,w(P)=w(P)$.
\end{proof}

\begin{corollary}\label{cor:proper}
Let $G\leq\operatorname{GL}(V)$ preserve $\beta$, let $a,b,e\in G$, and
let $K=\langle k_1,\ldots,k_m\rangle\leq G$.  Suppose that
$P\in k[a,b]$ is self-adjoint and
\[
  (s-I)P=0\qquad(s=a,b,e,k_1,\ldots,k_m).
\]
If some $g_0\in G$ satisfies $g_0w(P)\ne w(P)$, then every subgroup
$\langle a,b,c,d,e\rangle$ with $c\in K$ and $d\in C_G(a,b)$ is proper
in $G$.
\end{corollary}

\begin{proof}
Every element of $C_G(a,b)$ commutes with the algebra $k[a,b]$, so it
fixes $w(P)$ by Theorem~\ref{thm:canonical}.  The displayed equations
say that $a,b,e,k_1,\ldots,k_m$ fix $\im P$ pointwise.  They therefore
fix $w(P)$, which lies in $\im P$.  The same holds for every $c\in K$.
All five generators thus fix $w(P)$, whereas $g_0$ does not.
\end{proof}

The representation of an original group may be nonfaithful: if its
proposed generators generated that group, their images would generate
its image.  A moved-vector obstruction in the image therefore suffices.
The original group membership, centralizer and representation statements
remain necessary premises.  In particular, an unidentified semilinear
automorphism is not covered by this argument.

\paragraph{Exact matrix formula.}
Write $\beta(x,y)=x^{\mathsf T}Jy$, with $J$ invertible and symmetric.
Self-adjointness is $P^{\mathsf T}J=JP$.  The matrix $JP$ is then
symmetric, so its off-diagonal terms cancel in the quadratic expression.
If $d(P)=\operatorname{diag}(JP)$ is the column of its diagonal entries,
then
\[
 \beta(Px,x)=\sum_i d(P)_i x_i^2,\qquad
                  w(P)=J^{-1}\sqrt{d(P)}.
\]
The square root is taken entry by entry.  In particular, $w(P)\ne0$ if
and only if $d(P)\ne0$.  A self-adjoint nonzero $P$ need not have a
nonzero canonical vector: the diagonal condition must be checked.

\begin{proposition}\label{prop:linear}
Let $\mathcal A=k[a,b]$, and let $W$ be the subspace of $\mathcal A$
defined by self-adjointness and the fixed-image equations in
Corollary~\ref{cor:proper}.  For a specified $g\in G$, define
\[
 T_g(P)=(g-I)^{[2]}(J^{-1})^{[2]}\operatorname{diag}(JP),
\]
where the brackets mean entrywise squaring, not a matrix square.  Then
$T_g:W\longrightarrow V$ is $k$-linear and
\[
                 T_g(P)\ne0\quad\Longleftrightarrow\quad
                 g\,w(P)\ne w(P).
\]
If $P_1,\ldots,P_r$ is any basis of $W$, a certificate moved by this $g$
exists in $W$ if and only if $T_g(P_i)\ne0$ for some $i$.
\end{proposition}

\begin{proof}
The form of $W$ is linear: self-adjointness and each equation
$(s-I)P=0$ are linear in $P$.  The map $P\mapsto\operatorname{diag}(JP)$
is also linear, and the other two matrices defining $T_g$ are fixed.
Thus $T_g$ is linear.  Entrywise squaring commutes with matrix-vector
multiplication in characteristic two.  The displayed formula for $w(P)$
therefore gives
\[
                   T_g(P)=\bigl((g-I)w(P)\bigr)^{[2]}.
\]
Frobenius is injective, proving the equivalence.  Finally, a linear map
vanishes on $W$ exactly when it vanishes on a basis of $W$.
\end{proof}

\paragraph{Certificate and search.}
The algebra $k[a,b]$ has dimension at most $(\dim V)^2$.  One may build
it by starting with $I$ and repeatedly adjoining right products of a
current basis with $a$ and $b$ until their span is stable.  Every basis
element then has an expression in $a,b$, and stability proves that this
span is the entire algebra.  The equations defining $W$ form a finite
linear system.  Proposition~\ref{prop:linear} reduces the moved-vector
test to linear maps on its solution space; no rank-one search is needed.
If a generating set for the image of $G$ is supplied, testing all its
generators also decides whether this particular space $W$ contains a
canonical vector moved by $G$.

An accepting certificate needs only one expression for $P\in k[a,b]$,
the exact form and group representation premises, self-adjointness, the
fixed-image equalities for a complete generating set of $K$, and an actual
$g_0$ moving the vector calculated by the displayed formula.  It need not
certify completeness of the search algebra or of the solution space.
By contrast, a claim that this search finds no such witness requires
their complete bases; even that claim does not establish generation or
settle the original group classification.

\end{document}
